% !TeX spellcheck = en_US
\documentclass[aspectratio=1610]{beamer}
\usetheme{Madrid}
\setbeamercovered{transparent}
\setbeamertemplate{navigation symbols}{}
\setbeamertemplate{itemize items}[circle]
\setbeamertemplate{enumerate items}[default]

\makeatletter
\define@key{beamerframe}{wide}[10pt]{%
	\def\beamer@cramped{\itemsep #1\topsep0.5pt\relax}}
\makeatother
	

\usepackage[english]{babel}
\usepackage[utf8]{inputenc}
\usepackage[T1]{fontenc}
\usepackage{array}
\usepackage{xcolor}
\usepackage{colortbl}
\usepackage{booktabs}

\title[Winter - Lecture 3b] % (optional, use only with long paper titles)
{Intermediate Macroeconomics: Lecture  W3b }
\subtitle{ECON 311 -- Intermediate Macroeconomics}

\author{Muhebullah Karimzada}
\institute{School of Economics \\\ University of Northern British Columbia}
\date{Fall 2026}


% ---------- UNBC COLORS ----------
\usepackage{xcolor}
\definecolor{UNBCGreen}{RGB}{0,102,51}
\definecolor{UNBCDark}{RGB}{0,74,37}
\definecolor{UNBCGold}{RGB}{189,155,96}
\definecolor{SoftCream}{RGB}{249,247,242}
\definecolor{SoftGray}{RGB}{244,246,248}
\definecolor{AccentBlue}{RGB}{41,98,255}
\definecolor{AccentOrange}{RGB}{230,126,34}
\definecolor{AccentTeal}{RGB}{0,150,136}
\definecolor{AccentRed}{RGB}{192,57,43}
\setbeamercolor{normal text}{fg=black,bg=white}
\setbeamercolor{structure}{fg=UNBCDark}
\setbeamercolor{title}{fg=white,bg=UNBCDark}
\setbeamercolor{subtitle}{fg=UNBCGold}
\setbeamercolor{frametitle}{fg=white,bg=UNBCDark}
\setbeamercolor{block title}{fg=white,bg=UNBCDark}
\setbeamercolor{block body}{fg=black,bg=SoftGray}
\setbeamercolor{block title alerted}{fg=white,bg=AccentRed}
\setbeamercolor{block body alerted}{fg=black,bg=SoftGray}
\setbeamercolor{item}{fg=UNBCDark}
\setbeamercolor{section in toc}{fg=UNBCDark}
\setbeamercolor{alerted text}{fg=AccentRed}
\setbeamercolor{author in head/foot}{fg=white,bg=UNBCDark}
\setbeamercolor{date in head/foot}{fg=white,bg=UNBCDark}
% ---------- UNBC BRANDING ----------
\usepackage{graphicx}
\usepackage{lmodern}
\usepackage{tcolorbox}
\usepackage{tikz}
\usetikzlibrary{positioning,calc}
\setbeamertemplate{headline}{}
\setbeamertemplate{footline}{%
  \leavevmode%
  \hbox{%
  \begin{beamercolorbox}[wd=.78\paperwidth,ht=2.8ex,dp=1.2ex,leftskip=1em]{author in head/foot}%
  \color{white}\textbf{ECON 311} \hspace{0.5em} \textcolor{UNBCGold}{|} \hspace{0.5em} Intermediate Macroeconomics%
  \end{beamercolorbox}%
  \begin{beamercolorbox}[wd=.22\paperwidth,ht=2.8ex,dp=1.2ex,rightskip=1em plus1fil]{date in head/foot}%
  \color{white}\insertframenumber/\inserttotalframenumber%
  \end{beamercolorbox}}%
}

\newtcolorbox{mybox}[2][]{
  colback=#2!8,
  colframe=#2,
  boxrule=0.8pt,
  arc=2mm,
  left=2mm,right=2mm,top=1.5mm,bottom=1.5mm,
  #1
}

\newcommand{\topbanner}[1]{%
\begin{tikzpicture}[remember picture,overlay]
\fill[UNBCGold] ([yshift=-0.65cm]current page.north west) rectangle ([yshift=-0.48cm]current page.north east);
\node[anchor=north west,font=\small\bfseries,text=UNBCDark] at ([xshift=0.35cm,yshift=-0.78cm]current page.north west) {#1};
\end{tikzpicture}%
}
\begin{document}

\begin{frame}[plain]
  \titlepage
\end{frame}

\begin{frame}[wide]{What are You Going to Learn }
	\begin{itemize}
		\item Proximate the output trends/potential output
		\item Calculate moving average
		\item Introduce Hamilton Filter
	\end{itemize}
\end{frame}

\begin{frame}[wide]{Moving Averages (MA)}
	\begin{itemize}
		\item Calculating the average of fixed number of the most recent observations in a rolling window
		\vspace{2mm}
		\begin{itemize}
			\item e.g., MA of 3 quarters for 2001Q1 can be the average of 2000Q3, 2000Q4 and 2001Q1
			\item Then, we move to 2001Q2, the MA includes 2000Q4, 2001Q1 and 2001Q2
			\item Continue like this for the whole data series
		\end{itemize}
		\item The goal is to smooths out short-term fluctuations
		\item Highlights longer-term trends or cycles
	\end{itemize}
\end{frame}

\begin{frame}[wide]{Trailing and Centered Moving Averages}
	\begin{itemize}
		\item \textit{Trailing MA (Causal MA)}: Uses only past data points up to the current time.
		\begin{itemize}
			\item Example: A 5-period trailing MA at time \( t \) is:  
		\[
		MA_t = \frac{1}{5} \sum_{i=t-4}^{t} X_i
		\]
		\end{itemize}
		\item \textit{Centered MA (Symmetric MA)}: Averages data from both past and future points.
		\begin{itemize}
			\item Example: A 5-period centered MA at time \( t \) is:
		\[
		MA_t = \frac{1}{5} \sum_{i=t-2}^{t+2} X_i
		\]
		\end{itemize}
	\end{itemize}
\end{frame}

\begin{frame}[wide]{Leading Moving Average}
	\begin{itemize}
		\item \textit{Leading MA (Forward MA)}: Uses only future data points, which is uncommon in real-world applications.
		\vspace{1mm}
		\begin{itemize}
			\item Example: A 5-period leading MA at time \( t \) is:
		\[
		MA_t = \frac{1}{5} \sum_{i=t}^{t+4} X_i
		\]
		\end{itemize}
		\vspace{3mm}
		\item Most macroeconomics time series data would use Trailing MA, because economic agent at a specific time would not have any exact information about the future 
	\end{itemize}
\end{frame}

\begin{frame}[wide]{Simple Moving Average (SMA)}
	\begin{itemize}
		\item Basic form of moving average
		\item This is what we have been showing so far, for Trailing MA, 
		\[
		SMA_t = \frac{1}{N} \sum_{i=t-N+1}^{t} X_i
		\]
		\item Equal weights for all observations (the weight is $1/N$)
	\end{itemize}
\end{frame}

\begin{frame}[wide]{Simple Moving Average (SMA): Example}
	\begin{itemize}
		\item Suppose we have stock prices for the past 5 days:
	\end{itemize}
	\[
	X_{1} = 100, \quad X_{2} = 102, \quad X_{3} = 104, \quad X_{4} = 103, \quad X_5 = 105
	\]
	\begin{itemize}
		\item The 3-day SMA at time \( 3 \) is:
		\[
		SMA_{3} = \frac{100 + 102 + 104 }{3} = 102
		\]
		\item The 3-day SMA at time \( 4 \) is:
		\[
		SMA_4 = \frac{102 + 104 + 103 }{3} = 103
		\]
		\item The 3-day SMA at time \( 5 \) is:
		\[
		SMA_5 = \frac{104 + 103 + 105}{3} = 104
		\]
	\end{itemize}
\end{frame}

\begin{frame}[wide]{Weighted Moving Average (WMA)}
	\begin{itemize}
		\item Assigns different weights to observations (e.g., giving more importance to recent data)
		\item Unlike the Simple Moving Average (SMA), where all points are equally weighted, WMA emphasizes more recent values
		\item The formula for WMA over a window size \( N \) is:
		\[
		WMA_t = \frac{\sum_{i=t-N+1}^{t} w_{i-t} X_i}{\sum_{i=t-N+1}^{t} w_{i-t}}
		\]
		where \( w_{i-t} \) represents the assigned weight for each observation (e.g., $w_{-1}$ is the weight on the $X$ that is one period lag, $w_{-2}$ is the weight on the variable that is 2 period lag)
	\end{itemize}
\end{frame}

\begin{frame}[wide]{Drawbacks of Moving Averages}
	\begin{itemize}
		\item \textit{End-Point Problems}: 
		\vspace{1mm}
		\begin{itemize}
			\item At the start and end of the data series, fewer past or future values are available, leading to missing or biased estimates
		\end{itemize}
		\item \textit{Lag in Identifying Turning Points}: 
		\vspace{1mm}
		\begin{itemize}
			\item Since moving averages rely on past data, they respond slowly to sudden changes, delaying trend identification
		\end{itemize}
		\item 
		\textit{May Miss Structural Breaks}: 
		\vspace{1mm}
		\begin{itemize}
			\item Moving averages smooth out fluctuations, which may cause important shifts, such as economic crises or policy changes, to be overlooked
		\end{itemize}
		
	\end{itemize}
\end{frame}

\begin{frame}[wide]{Drawbacks of Moving Averages}
	\begin{itemize}
		\item \textit{Arbitrary Selection of Weights}:
		\vspace{1mm}
		\begin{itemize}
			\item In weighted moving averages, the choice of weights is subjective and may not always reflect true data dynamics
		\end{itemize}
		\item \textit{Sensitivity to Window Length Selection}: 
		\vspace{1mm}
		\begin{itemize}
			\item A short window reacts quickly but may capture too much noise, while a long window smooths better but introduces more lag
		\end{itemize}
	\end{itemize}
\end{frame}

\begin{frame}[wide]{Introduction to Hamilton Filter}
	\begin{itemize}
		\item Developed by James Hamilton from UCSD (Hamilton, 2018)
		\item The original paper shows it is an alternative to another popular smoothing method Hodrick-Prescott filter (we won't go into detail on the HP filter)
		\item Based on regression framework which provide more basic statistical properties
		\item Addresses traditional filter limitations
	\end{itemize}
\end{frame}

\begin{frame}[wide]{Hamilton Filter: Basic Concept}
	\begin{itemize}
		\item Projects future values based on recent past
		\item Focuses on $h$-period ahead prediction
		\item Uses linear regression methodology (more specifically, time series econometric)
		%\item Calculates trend deviations systematically
	\end{itemize}
\end{frame}

\begin{frame}[wide]{Hamilton Filter Mathematics}
	\begin{itemize}
		\item \textit{Basic Equation:} 
		\[
		y_{t+h} = \beta_0 + \beta_1 y_t + \beta_2 y_{t-1} + \beta_3 y_{t-2} + \beta_4 y_{t-3} + v_{t+h}
		\]
		The equation uses lagged values of \( y \) to predict \( y_{t+h} \), where \( v_{t+h} \) captures the cycle component (deviations from trend)
		
		\item \textit{Cycle Component:} 
		\[
		v_{t+h} = y_{t+h} - \text{predicted value}
		\]
		This isolates the cyclical component by subtracting the predicted trend value from the observed value
		
		\item \textit{Preserves Time Series Properties:} Unlike other filters, the Hamilton Filter does not require detrending or smoothing assumptions, retaining the original properties of the time series

	\end{itemize}
\end{frame}


\begin{frame}[wide]{Why 4 Lags and \( h = 8 \)?}
	\begin{itemize}
		\item \textit{Choice of 4 Lags:}
		\begin{itemize}
			\item Quarterly data often exhibit seasonal patterns 
			\item Using 4 lags captures one full year of quarterly dynamics, accounting for both seasonal effects
			\item Ensure the residual term is stationary -- i.e., bound up and down around zero (you will learn more in time series econometric)
		\end{itemize}
		
		\item \textit{Choice of \( h = 8 \):}
		\begin{itemize}
			\item \( h = 8 \) corresponds to a 2-year horizon for quarterly data, which is long enough to separate trends from short-term cycles
			\item This choice avoids blending cyclical movements with longer-term trends, ensuring clear distinction between the two

		\end{itemize}
	\end{itemize}
\end{frame}

\begin{frame}[wide]{Advantages of the Hamilton Filter}
	\begin{itemize}
		\item \textit{Theoretically Well-Founded:}
		\begin{itemize}
			\item Based on standard regression principles, it uses lagged observations to model trends and cycles.
			\item This avoids arbitrary smoothing parameters and is grounded in economic theory.
			\item Sometime refer as 
		\end{itemize}
		\vspace{4mm}
		\item \textit{More Reliable in Real-Time:}
		\begin{itemize}
			\item Because it uses only past data (causal filtering), it provides realistic estimates for real-time decision-making.
			\item No need for future data, making it ideal for forecasting and policy applications.
		\end{itemize}
		
	\end{itemize}
\end{frame}


\begin{frame}[wide]{Advantages of the Hamilton Filter}
	\begin{itemize}
		\item \textit{Preserves Basic Time Series Properties:}
		\begin{itemize}
			\item Does not distort key properties like stationarity, seasonality, or autocorrelation.
			\item Ensures that filtered data aligns closely with the original series.
		\end{itemize}
		\vspace{4mm}
		\item \textit{Robust to Different Data Frequencies:}
		\begin{itemize}
			\item Performs well with quarterly, monthly, or annual data, making it highly versatile.
			\item Recommended horizon \( h \) can be adjusted for different data frequencies (e.g., \( h = 8 \) for quarterly, \( h = 24 \) for monthly).
		\end{itemize}
		
		
	\end{itemize}
\end{frame}

\begin{frame}[wide]{Applying MA or Hamilton Filter}
	\begin{itemize}
		\item Recall the output decomposition formula:
		\vspace{4mm}
		\[
		Y_t = \bar{Y}_t + \bar{Y}_t \tilde{Y}_t
		\]
		\item The term \( \beta_0 + \beta_1 y_t + \beta_2 y_{t-1} + \beta_3 y_{t-2} + \beta_4 y_{t-3} \) can be interpreted as the potential output component (\( \bar{Y}_t \))
		\item The term \( v_{t+h} \) corresponds to the cyclical component (\( \bar{Y}_t \tilde{Y}_t \))
		\item To improve statistical properties and simplify analysis, GDP is often transformed by taking its logarithm
	\end{itemize}
\end{frame}





\end{document}